% Standalone problem document, generated from the adjacent README.md.
% Compile with XeLaTeX. Mathematical content is maintained in Markdown.
\documentclass[11pt,a4paper]{article}
\usepackage[fontsize=10.5pt]{scrextend}
\usepackage[margin=22mm,headheight=15pt,headsep=9mm,footskip=11mm]{geometry}
\usepackage{fontspec}
\setmainfont{texgyrepagella}[Extension=.otf,UprightFont=*-regular,BoldFont=*-bold,ItalicFont=*-italic,BoldItalicFont=*-bolditalic]
\setsansfont{texgyreheros}[Extension=.otf,UprightFont=*-regular,BoldFont=*-bold,ItalicFont=*-italic,BoldItalicFont=*-bolditalic]
\setmonofont{texgyrecursor}[Extension=.otf,UprightFont=*-regular,BoldFont=*-bold,ItalicFont=*-italic,BoldItalicFont=*-bolditalic,Scale=MatchLowercase]
\usepackage{amsmath,amssymb}
\usepackage{unicode-math}
\setmathfont{texgyrepagella-math.otf}
\usepackage{xcolor,microtype,enumitem,needspace,fancyhdr}
\usepackage{bookmark}
\definecolor{ink}{HTML}{17344B}
\definecolor{accent}{HTML}{197D80}
\definecolor{muted}{HTML}{596773}
\hypersetup{colorlinks=true,linkcolor=accent,urlcolor=accent,pdftitle={MF-14 - Degree
coverage with seven matrix
multiplications},pdfauthor={Open Problems in Numerical Linear Algebra}}
\urlstyle{same}
\setlength{\parindent}{0pt}
\setlength{\parskip}{5pt plus 1pt minus 1pt}
\linespread{1.04}
\setlength{\emergencystretch}{3em}
\widowpenalty=10000
\clubpenalty=10000
\displaywidowpenalty=10000
\allowdisplaybreaks[1]
\setlist{nosep,leftmargin=1.5em}
\providecommand{\tightlist}{\setlength{\itemsep}{0pt}\setlength{\parskip}{0pt}}
\providecommand{\pandocbounded}[1]{#1}
\pagestyle{fancy}
\fancyhf{}
\fancyhead[L]{\sffamily\footnotesize\color{muted}OPEN PROBLEMS IN NUMERICAL LINEAR ALGEBRA}
\fancyhead[R]{\sffamily\bfseries\footnotesize\color{accent}MF-14}
\fancyfoot[L]{\parbox[b]{0.9\textwidth}{\sffamily\fontsize{8}{10}\selectfont\color{muted}Editorial ratings; a bounded search cannot certify that a problem remains unsolved.\\Literature check: 2026-09-08}}
\fancyfoot[R]{\sffamily\footnotesize\color{muted}\thepage}
\renewcommand{\headrulewidth}{0.3pt}
\renewcommand{\headrule}{\hbox to\headwidth{\color{accent}\leaders\hrule height \headrulewidth\hfill}}
\makeatletter
\renewcommand{\section}{\Needspace{5\baselineskip}\@startsection{section}{1}{0pt}{12pt plus 2pt minus 2pt}{4pt}{\sffamily\bfseries\large\color{ink}}}
\renewcommand{\subsection}{\Needspace{4\baselineskip}\@startsection{subsection}{2}{0pt}{9pt plus 1pt minus 1pt}{3pt}{\sffamily\bfseries\normalsize\color{ink}}}
\makeatother
\setcounter{secnumdepth}{0}
\begin{document}
{\sffamily\small\color{muted}Matrix functions and stability\par}
\vspace{3pt}
{\raggedright\hyphenpenalty=10000\exhyphenpenalty=10000\sffamily\bfseries\fontsize{21}{24}\selectfont\color{ink}Degree
coverage with seven matrix multiplications\par}
\vspace{7pt}
{\sffamily\small\textbf{Difficulty:} challenging\quad\textbf{Importance:} interesting
to the community\par}
\vspace{4pt}
{\color{accent}\hrule height 0.8pt}
\vspace{6pt}
\subsection{Problem statement}\label{problem-statement}

Start with the scalar polynomials \(1,x\in\mathbb C[x]\). Allow
arbitrarily many complex linear combinations of already computed
polynomials and at most seven multiplications of two already computed
polynomials. Let \(\mathcal P_7\subseteq\mathbb C[x]_{\le128}\) be the
set of possible outputs.

Identify a polynomial with its coefficient vector in \(\mathbb C^{129}\)
and let \(\overline{\mathcal P_7}^{\,Z}\) be its Zariski closure: the
common zero set of all polynomial equations in the coefficients that
vanish throughout \(\mathcal P_7\). Is

\[\max\{d\in\{0,\ldots,128\}:\mathbb C[x]_{\le d}\subseteq\overline{\mathcal P_7}^{\,Z}\}=42?\]

The target concerns closure and the complex coefficient field. Exact
representability of every polynomial is a stronger assertion.

\subsection{Why it matters}\label{why-it-matters}

For a dense input matrix, matrix products dominate the cost of
polynomial evaluation. The conjecture determines how much polynomial
degree a budget of seven products can cover, providing a theoretical
benchmark for matrix-function evaluation schemes.

\subsection{References}\label{references}

\begin{enumerate}
\def\labelenumi{\arabic{enumi}.}
\tightlist
\item
  E. Jarlebring and G. Lorentzon, \emph{The polynomial set associated
  with a fixed number of matrix-matrix multiplications},
  arXiv:2504.01500v3 (13 August 2025), §2.2 (closure convention), §5.4,
  Conjecture 13. \href{https://arxiv.org/html/2504.01500}{Primary text}.
\item
  J. Sastre, J. Ibáñez, J. M. Alonso, and E. Defez, \emph{Beyond
  Paterson--Stockmeyer: Advancing matrix polynomial computation}, WSEAS
  Transactions on Mathematics 24 (2025), 684--693, §4 (a five-product
  construction).
  \href{https://wseas.com/journals/mathematics/2025/b385106-036\%282025\%29.pdf}{Primary
  paper}, \href{https://doi.org/10.37394/23206.2025.24.68}{DOI}.
\item
  J. M. Alonso, J. Sastre, J. Ibáñez, and E. Defez, \emph{A systematic
  framework for stable and cost-efficient matrix polynomial evaluation},
  arXiv:2603.23143v1 (24 March 2026), abstract and conclusions.
  \href{https://arxiv.org/html/2603.23143}{Primary text}.
\end{enumerate}

\subsection{Status check --- 2026-09-08}\label{status-check-2026-09-08}

The latest arXiv source remains v3 and labels the displayed equality a
conjecture supported by computations. The later papers give
constructions and stability procedures without establishing the
seven-product degree-coverage claim. Searches included the exact title
with \textit{2026},
\textit{Jarlebring Lorentzon conjecture seven multiplications 42},
and \textit{Sastre polynomial evaluation 2026}. No resolution was
located. A
\href{https://indico3.mpi-magdeburg.mpg.de/event/58/contributions/1074/}{3
September 2026 author talk} also describes higher-cost degree coverage
as conjectural. Six-product real/complex formulations have a separate
source ambiguity and are deliberately not included here.
\end{document}
